\documentclass[11pt]{article}
\usepackage[T1]{fontenc}
\usepackage{lmodern,amsmath,amssymb,amsthm,mathtools,booktabs,microtype}
\usepackage[margin=1in,top=0.87in,bottom=0.90in]{geometry}
\usepackage[table]{xcolor}
\usepackage{fancyhdr}
\usepackage[colorlinks=true,linkcolor=teal!55!black,citecolor=teal!55!black,urlcolor=teal!55!black]{hyperref}
\definecolor{ink}{HTML}{17344A}
\definecolor{accent}{HTML}{16665F}
\newtheorem{theorem}{Theorem}[section]
\newtheorem{proposition}[theorem]{Proposition}
\newtheorem{lemma}[theorem]{Lemma}
\newtheorem{corollary}[theorem]{Corollary}
\newcommand{\N}{\mathbb{N}}
\newcommand{\R}{\mathbb R}
\pagestyle{fancy}\fancyhf{}\fancyhead[L]{\footnotesize\sffamily Keio and Shaitetsu}\fancyhead[R]{\footnotesize\sffamily Diophantine multiplicities and zero distributions}
\fancyfoot[C]{\sffamily\footnotesize\thepage}\renewcommand{\headrulewidth}{0.25pt}
\setlength{\headheight}{14pt}\setlength{\parskip}{2pt}
\setlength{\emergencystretch}{1.8em}\raggedbottom
\hypersetup{pdftitle={Unbounded multiplicities in the Erdos--Graham equation},pdfauthor={Keio and Shaitetsu}}
\begin{document}
\thispagestyle{plain}
{\sffamily\small\color{accent} RESEARCH MANUSCRIPT \hfill 16 NOVEMBER 2026}\par
\vspace{10pt}
{\color{ink}\sffamily\bfseries\fontsize{24}{28}\selectfont Unbounded multiplicities in the\\ Erd\H{o}s--Graham equation\par}
\vspace{5pt}
{\large\sffamily Sharp corner transversals, prime-zeta asymptotics, and real-zero tests}\par
\vspace{9pt}{\normalsize Keio and Shaitetsu}\par
\vspace{7pt}\noindent\textcolor{accent}{\rule{\linewidth}{0.6pt}}
\medskip

For a fixed number of summands $k$, let $N(k)$ count the positive integer
solutions of
\[
 \frac{n}{2^n}=\sum_{i=1}^k\frac{a_i}{2^{a_i}},\qquad
 a_1<\cdots<a_k,\qquad k\ge2 .
\]
Tengely, Ulas and Zygad\l o proved in 2020 that $N(k)\ge5$ on an
arithmetic progression and conjectured that $\limsup N(k)=\infty$.
We prove their Conjecture~3.3. For every $r$ an entire tail of an
explicit progression consists of values with at least $r+1$ solutions.
On a subsequence,
\[
 \limsup_{k\to\infty}N(k)\sqrt{\frac{\log\log k}{\log k}}
 \ \ge\ \sqrt{\frac{2}{\log 2}} .
\]
The construction deletes one entry from a consecutive block and
synchronizes the resulting congruences. We classify every possible
single deletion and determine its exact mean contribution.

\medskip
\noindent\textbf{Companion results.} We settle two conjectures posed in the
Verifiable Mathematical Reasoning Corpus (VMRC). For diagonal geometric
weights, exactly two color classes minimize a directed corner
transversal. At the right zero edge of the prime zeta function, we
establish the Gaussian prefactor relative to the corpus zero-density
formula and determine conditional higher-harmonic deficits.
For the operator question of Csordas and Escassut we classify the
symbols on periodic strip functions and on a class of infinite pure
point spectra with endpoint atoms. A trace-form test counts nonreal
pairs in the periodic case. A final algebraic result classifies
compatibility in two families of decoder identities.

\medskip
\noindent\textbf{Scope.} The fixed-$k$ theorem is unconditional and
self-contained apart from the prime number theorem in arithmetic
progressions used for its quantitative refinement. The question whether
every prescribed $n$ has a representation and the unrestricted
Csordas--Escassut classification remain outside the conclusions.
The radial zero-density input for the prime-zeta application is
identified in Section~\ref{sec:primeedge}. This manuscript presents
proofs for review; historical priority has not been independently established.
\newpage

\section{Many solutions with a prescribed number of terms}
\label{sec:tuz}
Tengely, Ulas and Zygad\l o count different tuples
$(n;a_1,\ldots,a_k)$ for fixed $k$ \cite[Conjecture 3.3]{tuz}.
The left endpoint $n$ may vary. Their construction uses two displaced
final terms and a discrete logarithm. VMRC~\cite{vm090} encodes each
representation with a normalized integer remainder. For the present
purpose it is convenient to reverse that computation: specify the
positions that occur and recover $n$ from the terminal remainder.

\begin{lemma}[One omitted position]\label{lem:hole}
Put $d=k+1$. Select the offsets $1,\ldots,d$ except $j=d-t$,
where $1\le t\le d-1$. The selected exponents $a=n+i$ solve the
equation precisely when $n$ is a positive integer and
\begin{equation}
 (n+j)(1+2^t)=2^{k+2}-(t+2).\label{eq:onehole}
\end{equation}
For $t\ge4$ integrality is equivalent to the existence of an integer
$m\ge3$ such that $t=2^m-2$ and $k+2\equiv m\pmod{2t}$.
For $t=1$ integrality never holds. For $t=2$ and $3$ the required
residues of $k+2$ are respectively $2\bmod4$ and $5\bmod6$.
\end{lemma}
\begin{proof}
Writing $b_i=1$ for a selected offset and $b_i=0$ otherwise gives
the exact integer recursion
$R_0=n$, $R_i=2R_{i-1}-b_i(n+i)$.
The equation is equivalent to $R_d=0$. Equivalently, insert
$\sum_{i=1}^d(n+i)2^{-i}=n+2-(n+d+2)2^{-d}$,
subtract $(n+j)2^{-j}$, and multiply the equality by $2^d$.
This gives \eqref{eq:onehole}.

Modulo $2^t+1$ we have $2^t\equiv-1$. For $t\ge2$ the residues of the
first $2t$ powers of two are the $t$ numbers $2^b$ and the $t$
numbers $2^t+1-2^b$, $0\le b<t$. The two lists are disjoint
(the first ends at $2^{t-1}$ and the second starts above it);
each list has distinct entries, so the period is exactly $2t$.
For $t\ge4$ the smallest number in the second list, $2^{t-1}+1$,
exceeds $t+2$. Thus $t+2$ must itself be a power of two,
and the exponent in the congruence must be the stated residue.
Reduction modulo $3,5,9$ gives the three small cases.
\end{proof}

\begin{theorem}\label{thm:tuz}
For every $r\ge1$, there are $L_r\ge1$ and $b_r$ such that
$N(k)\ge r+1$ for every sufficiently large $k\equiv b_r\pmod{L_r}$.
The $r+1$ solutions can be chosen with different values of $n$.
\end{theorem}
\begin{proof}
Choose different primes $p_1,\ldots,p_r\equiv1\pmod4$.
For each one put
\[
 A_p=2^p-1,\quad t_p=2A_p=2^{p+1}-2,\quad m_p=p+1 .
\]
We impose $k+2\equiv p+1\pmod{4A_p}$ for all of them. The odd
numbers $A_p$ are pairwise coprime, by
$\gcd(2^p-1,2^q-1)=2^{\gcd(p,q)}-1$.
All target residues are $2\bmod4$. Writing $k+2=2+4h$ therefore
reduces the problem to the ordinary Chinese remainder theorem for
$h\equiv(p-1)/4\pmod{A_p}$. The resulting period is
$L_r=4\prod_p A_p$.

Each omitted position is internal for large $k$. Equation
\eqref{eq:onehole} gives $n=2^{k+2}/(1+2^{t_p})+O_r(k)$.
The positive leading coefficients are distinct. An additional solution
uses every offset $1,\ldots,k$, with
$n=2^{k+1}-k-2$, as follows by the same geometric sum.
Its leading coefficient relative to $2^{k+2}$ is $1/2$,
different from all the others. This proves the assertion.
\end{proof}

For an alternative elementary proof of the qualitative conclusion, start
with $m_1=2$. Having chosen $m_1,\ldots,m_s$, write $t_i=2^{m_i}-2$ and
set $m_{s+1}=m_s+\operatorname{lcm}(2t_1,\ldots,2t_s)$.
Induction gives $m_s\equiv m_i\pmod{2t_i}$ for every $i<s$.
The progression $k+2\equiv m_r\pmod{\operatorname{lcm}(2t_1,\ldots,2t_r)}$
works for all $r$ omissions. This proof of Conjecture~3.3 needs
only finite geometric sums and elementary congruences.

\begin{corollary}\label{cor:rate}
With natural logarithms,
$\displaystyle\limsup_{k\to\infty}N(k)\sqrt{\log\log k/\log k}
\ge\sqrt{2/\log2}$.
\end{corollary}
\begin{proof}
In Theorem~\ref{thm:tuz} use the first $r$ primes congruent to $1$
modulo four, and take $k\in[L_r,2L_r)$ in the required class.
For each selected $t$, $t\le L_r/2\le k/2$; the right-hand quotient
in \eqref{eq:onehole} exceeds $2^{k-t}$ for $k\ge6$ and these
$t\ge2$, so $n>0$. Different omissions determine different normalized
offset sets and therefore different solutions even without a size
threshold depending on $r$. By the prime number theorem in arithmetic
progressions $p_i\sim2i\log i$. Consequently
$\sum_{i\le r}p_i\sim r^2\log r$, and
$\log k=\log L_r+O(1)\sim(\log2)r^2\log r$.
Substitution of $N(k)\ge r+1$ gives the constant.
\end{proof}

\begin{proposition}[Mean contribution of single omissions]
Let $H(k)$ count only solutions of the internal one-omission form
in Lemma~\ref{lem:hole}, with the consecutive solution excluded. Then
\[
 \lim_{X\to\infty}\frac1X\sum_{2\le k\le X}H(k)
   =\frac16+\frac14\sum_{j\ge1}\frac1{2^j-1}.
\]
\end{proposition}
\begin{proof}
The exceptional offsets $t=2,3$ contribute respective densities $1/4$
and $1/6$. For $m\ge3$, $t=2^m-2$ contributes $1/(2t)$.
In a single class, discarding $k<2t$ and the finitely many small
$k<6$ changes the count by $O(1)$ uniformly in $t$, and all later
parameters satisfy $n>0$ as above. For $k\le X$ only $O(\log X)$
possible offsets occur, since $t\le k$. The total error on summing
their progression counts is $O(\log X)$, and the omitted tail of
the sum of densities tends to zero. Substituting $j=m-1$ gives the
displayed limit.
\end{proof}

\noindent\textbf{Check at ordinary sizes.}
For $k=8,t=2$ the deletion is $j=7$ and $n=197$. The eight offsets
are $1,2,3,4,5,6,8,9$. At $k=128$ both $t=2$ and $t=62$ work:
$130\equiv2\pmod4$ and $130\equiv6\pmod{124}$.
For $t=62$ the formula gives $n=295147905179352825725$.
The source's Question~3.4 asks whether its Conjecture~3.2 implies
Conjecture~3.3. The consequent now follows unconditionally. Their
separate infinitude conjecture for offsets modulo $2^{u+3}-3$ uses
a different pattern; the present construction does not establish it.

\subsection{Density of all multiplicity thresholds for one omission}

A fixed finite set of missing offsets also admits a finite-state
congruence test. Take $F\subset\N$ nonempty and fixed, $s=|F|$,
$d=k+s$, and delete $j=d-t$ for every $t\in F$ whenever $d>\max F$.
Eliminating the normalized remainder gives
\[
 Q_F(n+d)=2^{d+1}-C_F,\qquad
 Q_F=1+\sum_{t\in F}2^t,\qquad C_F=2-\sum_{t\in F}t2^t .
\]
Indeed each deleted position costs $(n+d-t)2^t$ in the final integer
remainder equation. The odd modulus $Q_F$ makes
$z\mapsto 2z\bmod Q_F$ a permutation. Starting at $1$, its exact
period is $\operatorname{ord}_{Q_F}(2)$. Thus, for all large $k$,
this template occurs on a single residue class or never occurs, according
as $C_F$ belongs to the orbit or not. Several templates can be tracked
simultaneously by the product permutation. This is a finite residue
version of digit propagation; it avoids searching through unknown $n$.

\begin{theorem}[Limiting distribution of the one-omission count]
For each integer $r\ge1$ the set of $k$ satisfying $H(k)\ge r$
has a natural density $d_r$. It satisfies
\[
 0<d_r\le \sum_{m\ge r}\frac{1}{2(2^m-2)}
 \quad (r\ge3).
\]
In particular the one-omission counts have a limiting distribution
of unbounded support and exponentially decreasing upper tail.
\end{theorem}
\begin{proof}
Track $t=3$ and the templates $t=2^m-2$, $2\le m<J$.
After finitely many initial exceptions, their count is periodic by the
finite-state description. For $k\le X$, only $m\le\log_2(X+2)$ can
occur. The upper density of the set on which the truncated count differs
from $H(k)$ is at most $\sum_{m\ge J}1/[2(2^m-2)]$: a class of period
$2t$ contributes at most $X/(2t)+1$ terms and the sum of the unit errors
is $O(\log X)$. The bound tends to zero as $J\to\infty$, giving existence
of the threshold densities. With $J=r$, at most $r-1$ initial templates
are tracked, which yields the displayed upper bound. Theorem~\ref{thm:tuz}
constructs a progression of positive density with $r$ genuine omissions.
Finally $d_r-d_{r+1}$ are the point probabilities for $r\ge1$, and
$1-d_1$ is the mass at zero. Positivity for every threshold and
$d_r\to0$ give unbounded support.
\end{proof}

For reproducibility the following Python fragment checks every admissible
single omission for $2\le k\le256$ against the original equation after
clearing the power-of-two denominator. It counts 143 positive tuples.
The general proof uses no numerical cutoff.
\begin{verbatim}
count = 0
for k in range(2, 257):
    d = k + 1
    for t in range(1, d):
        j = d - t
        q, r = divmod(2**(k+2) - t - 2, 2**t + 1)
        m = (t+2).bit_length() - 1
        good = (t == 3 and (k+2)%6 == 5) or (
            2**m == t+2 and (k+2-m)%(2*t) == 0)
        assert (r == 0) == good
        n = q - j
        if r == 0 and n > 0:
            assert n*2**d == sum((n+i)*2**(d-i)
                for i in range(1,d+1) if i != j)
            count += 1
assert count == 143
\end{verbatim}

\section{Diagonal corner transversals and their unique extremizers}
\label{sec:diagonal-corners}

Put $T_d=\{(i,j)\in\mathbb Z_{\ge0}^2:i+j\le d\}$ and
$H_{i,j}=\{(i,j),(i+1,j),(i,j+1)\}$. A transversal is a set meeting
every $H_{i,j}$. For $s\in\{0,1,2\}$ let
\[
 C_s=\{(i,j):i+2j\equiv s\pmod3\}.
\]
Each $C_s$ meets every corner in exactly one vertex.

\begin{theorem}\label{thm:diagonal-weight}
Suppose $w_0\ge w_1\ge\cdots\ge0$ and
$\sum_{d\ge0}(d+1)w_d<\infty$. For every transversal $S$,
\[
 \sum_{(i,j)\in S}w_{i+j}
 \ \ge\
 \sum_{d\ge0}(w_d-w_{d+1})
       \left\lfloor\frac{(d+1)(d+2)}6\right\rfloor.
\]
Both $C_1$ and $C_2$ attain equality. If $w_d>w_{d+1}$ for all $d$,
these are the only minimizers. More generally, equality holds precisely
when $|S\cap T_d|=\lfloor(d+1)(d+2)/6\rfloor$ for every $d$ at which
$w_d>w_{d+1}$.
\end{theorem}

\begin{proof}
Write $b_d=\lfloor(d+1)(d+2)/6\rfloor$ and $b_{-1}=0$.
We first recall the finite monomial transversal theorem of Geramita,
Gregory and Roberts \cite[Theorem~5.4(1)]{ggr}. Among the degree-$d$
monomials in three variables, the minimum number needed to contain a
multiple of every degree-$(d-1)$ monomial is
$\lfloor\binom{d+2}{2}/3\rfloor=b_d$. The correspondence
$(i,j)\mapsto X^iY^jZ^{d-i-j}$ identifies the three vertices of each
corner contained in $T_d$ with the three degree-$d$ multiples of the
associated degree-$(d-1)$ monomial. Consequently $|S\cap T_d|\ge b_d$;
the same bound at $d=0$ is immediate.

Put $s_d=b_d-b_{d-1}$. Direct subtraction gives
\[
 s_{3r}=r,\qquad s_{3r+1}=s_{3r+2}=r+1\qquad(r\ge0).
\]
On the diagonal $i+j=d$, the congruence defining $C_s$ is
$i\equiv2d-s\pmod3$. For $s=1,2$ it gives exactly $s_d$ cells,
so both cosets attain all finite lower bounds simultaneously.
Since $w_d\to0$, nonnegative summation gives the exact identity
\[
 \sum_{(i,j)\in S}w_{i+j}
 =\sum_{d\ge0}(w_d-w_{d+1})|S\cap T_d|.
\]
This proves the value and the stated general equality criterion.

For strict decrease all finite bounds are equalities. To classify their
simultaneous equality cases, let
$S_n=\{i:(i,n-i)\in S\}\subseteq\{0,\ldots,n\}$, so $|S_n|=s_n$.
Each unmarked lower index $i\notin S_n$ requires $i$ or $i+1$ in
$S_{n+1}$. One upper mark can meet at most two such requirements, and
then their indices are consecutive. When $n=3r+2$, the $2r+2$
unmarked lower indices exactly exhaust the capacity of the $r+1$ upper
marks. Therefore each consecutive run of unmarked indices has even
length. In a run $[h,h+2l-1]$ the required upper marks are
$h+1,h+3,\ldots,h+2l-1$; this follows by successively pairing its
leftmost two indices.

We have $S_0=\varnothing$ and $S_1$ is either $\{0\}$ or $\{1\}$.
Interchanging $i$ and $j$ exchanges these cases. Suppose
$S_1=\{0\}$. Then $S_2$ is either $\{1\}$ or $\{2\}$. The first
violates the even-run condition, so $S_2=\{2\}$ and $S_3=\{1\}$.

Inductively suppose $r\ge1$ and
$S_{3r}=\{1,4,\ldots,3r-2\}$. The unmarked runs on this layer
are $[0,0]$ and the $r$ pairs $[3h-1,3h]$, $1\le h\le r$.
Since there are precisely $r+1$ upper marks, we must have
\[
 S_{3r+1}=\{e,3,6,\ldots,3r\},\qquad e\in\{0,1\}.
\]
If $e=1$, the next requirements come from the three singletons
$0,2,3r+1$ and the $r-1$ pairs between successive marked indices
$3,6,\ldots,3r$. Distinct runs require distinct upper marks, giving
at least $r+2$, impossible. Hence $e=0$. The $r$ pairs
$[3h-2,3h-1]$, $1\le h\le r$, on this layer now force
$2,5,\ldots,3r-1$ on the next layer, with one extra mark to cover
$3r+1$. It is either $3r+1$ or $3r+2$. The first leaves a terminal
unmarked singleton in $S_{3r+2}$, which again violates the even-run
condition. Thus $S_{3r+2}=\{2,5,\ldots,3r+2\}$. Pairing its
unmarked runs forces
$S_{3r+3}=\{1,4,\ldots,3r+1\}$, completing the induction.
These layers form $C_2$; reflection gives exactly $C_1$.
\end{proof}

\begin{corollary}\label{cor:diagonal-corners}
For every real $a>1$, the minimum diagonal weight in the corner problem
of \cite{vm140} is
\[
 \kappa(a,a)=\left(\sum_{i,j\ge0}a^{-i-j}\right)
                 \frac{a}{a^2+a+1}.
\]
The minimizers are exactly $C_1$ and $C_2$.
\end{corollary}
\begin{proof}
Set $q=1/a$. The generating function of the shell sizes of $C_1$ or
$C_2$ is
\[
 \sum_{d\ge0}s_d q^d
 =\sum_{r\ge0}\bigl(rq^{3r}+(r+1)(q^{3r+1}+q^{3r+2})\bigr)
 =\frac{q}{(1-q)(1-q^3)}.
\]
Division by total weight $(1-q)^{-2}$ gives
$q/(1+q+q^2)=a/(a^2+a+1)$. All successive weight differences are
positive, so Theorem~\ref{thm:diagonal-weight} gives uniqueness.
\end{proof}

\section{The Gaussian prefactor at the prime-zeta zero edge}
\label{sec:primeedge}

Let $\sigma_e>1$ be the unique root of
$\sum_{p\ge3}p^{-\sigma_e}=2^{-\sigma_e}$, where $p$ runs over primes.
Write $a_p=p^{-\sigma}$, $x=2^{-\sigma}$, and let
$S=\sum_{p\ge3}a_pU_p$ for independent Haar-uniform unit complex
numbers $U_p$. For $1<\sigma<\sigma_e$ put
\[
 L(z)=\sum_{p\ge3}\log I_0(a_pz),\qquad
 L'(\lambda)=x,\qquad I=\lambda x-L(\lambda).
\]
Here $I_0$ is the modified Bessel function; the real logarithm defines
$L$ on the positive real axis and its local analytic continuation is
specified below. There is a unique positive $\lambda$, since $L'(0)=0$,
$L''>0$, and $L'(\lambda)\uparrow\sum a_p>x$. Moreover
$\lambda\to\infty$ as $\sigma\uparrow\sigma_e$, by dominated convergence
and strict inequality at every finite $\lambda$ at the edge. Define
\[
 D_1(s)=(\log2)2^{-s}-\sum_{p\ge3}(\log p)p^{-s}.
\]
In particular $D_1(\sigma_e)=-\sum_{p\ge3}\log(p/2)p^{-\sigma_e}<0$.

\begin{theorem}[Gaussian prefactor]
Let $\rho(\sigma)$ be the horizontal zero-intensity density per unit vertical height of the
prime zeta function on $\Re s>1$, with the normalization of
\cite{vm004}. Under the tilted probability law
$d\mathbb P_\lambda=e^{\lambda\Re S-L(\lambda)}d\mathbb P$, let
$v_1=\operatorname{Var}_\lambda(\Re S)$ and
$v_2=\operatorname{Var}_\lambda(\Im S)$. Then
\[
 \rho(\sigma)=
 \frac{D_1(\sigma_e)^2}{2\pi\sqrt{v_1v_2}}e^{-I}(1+o(1))
 \qquad(\sigma\uparrow\sigma_e).
\]
\end{theorem}

\begin{proof}
We use the radial reduction (10.7) of \cite{vm004}. In the present
notation it gives $\rho(\sigma)=D_1(\sigma_e)^2
(1+O(\sqrt{\epsilon})) f_S(x,0)$ with
$\epsilon=\sum a_p-x\to0$ and $f_S$ the planar density in the original
law. The density under the tilt is denoted by $g$. Exact change of
measure gives $f_S(x,0)=e^{-I}g(x,0)$. For clarity, on the aligned zero fiber $S=x$ one has $U_2=-1$ and $\sum a_p|1-U_p|^2=2\epsilon$. The derivative at the zero differs in modulus from $|D_1(\sigma)|$ by at most $(2\epsilon\sum a_p(\log p)^2)^{1/2}$, by Cauchy--Schwarz. The sum is bounded uniformly near the edge. Averaging over the angle of the first prime preserves the value of the radial density at $(x,0)$. This explains the normalization of the complex Kac--Rice weight in the cited identity. It suffices to prove a local
central limit estimate for $g$ at its mean.

For real $u,v$ near zero, the centered characteristic function in the
tilted law is
\begin{equation}
 \phi(u,v)=\exp\{L(w)-L(\lambda)-iux\},\qquad
 w=\sqrt{(\lambda+iu)^2-v^2},
 \label{eq:edge-cf}
\end{equation}
where the branch takes value $\lambda$ at the origin. Indeed the angular
integral of $\exp(A\cos\theta+B\sin\theta)$ against normalized angular
measure is $I_0(\sqrt{A^2+B^2})$, first by real rotations for $A,B$
real and then by analytic continuation. Independence gives
\eqref{eq:edge-cf}. Rotation invariance before tilting gives the
joint cumulant function $L(\sqrt{s^2+t^2})$ for real arguments.
Differentiating at $(s,t)=(\lambda,0)$ yields
\[
 v_1=L''(\lambda),\qquad v_2=L'(\lambda)/\lambda=x/\lambda,
 \qquad \operatorname{Cov}_\lambda(\Re S,\Im S)=0.
\]

\emph{Variance and analytic bounds.}
Constants below are uniform for $\sigma$ in a fixed compact
neighborhood of $\sigma_e$ contained in $(1,\infty)$, with its
intersection with $(1,\sigma_e)$ understood. Take
$P=\lambda^{1/\sigma}$ and $M=\#\{p:P\le p\le2P\}
\asymp P/\log P$. Put $h=\log I_0$. There are $c,C>0$ such that on
$|z-\lambda|\le c\lambda$,
\[
 |L'''(z)|\le C M/\lambda^3,\qquad
 cM/\lambda^2\le L''(\lambda)\le C M/\lambda^2 .
\]
For details, $h$ has an analytic branch in a narrow cone about the
positive real axis. On bounded pieces this follows from $I_0(t)>0$
for $t\ge0$, compactness and $I_0(0)=1$. At infinity the sectorial
expansion $I_0(z)=e^z(2\pi z)^{-1/2}(1+O(1/z))$ applies
\cite[\S10.40]{dlmf}. Cauchy's estimates on a narrower cone give
$|h^{(j)}(z)|\ll(1+|z|)^{-j}$ for $j=2,3$.
The small-argument expansion $h(z)=O(z^2)$ gives local normal
convergence of $L$. In the differentiated sums each prime $p\le P$
costs $O(\lambda^{-j})$, and the remaining terms cost at most
$O(\sum_{p>P}p^{-j\sigma})=O(M\lambda^{-j})$ by the prime number
theorem and dyadic summation. For the lower bound the primes
$P\le p\le2P$ have $a_p\lambda$ in a fixed compact subinterval of
$(0,\infty)$. On it $h''$ is bounded below by a positive constant,
since it is the variance of $\cos\theta$ in the corresponding
strictly positive angular density.

\emph{Central limit.}
Fix $s,t\in\mathbb R$ and put $u=s/\sqrt{v_1}$ and $v=t/\sqrt{v_2}$.
Then $u=O(\lambda/\sqrt M)$ and $v=O(\sqrt\lambda)$. The exact
identity $w-(\lambda+iu)=-v^2/(w+\lambda+iu)$ gives
\[
 w=\lambda+iu-\frac{v^2}{2\lambda}+o(1).
\]
Indeed the error is $O(v^2|u|/\lambda^2+v^4/\lambda^3)=o(1)$.
Taylor expansion of $L$ gives
\[
 L(w)-L(\lambda)-iux=-\tfrac12v_2v^2-\tfrac12v_1u^2+o(1).
\]
The cubic remainder is $O(M\lambda^{-3}(|u|+1)^3)=o(1)$.
The error from the quadratic cross terms is
$O(v_1(|u|+1))=o(1)$. Thus
$\phi(s/\sqrt{v_1},t/\sqrt{v_2})\to e^{-(s^2+t^2)/2}$.

\emph{Integrable Fourier tails.}
Fix a sufficiently small $\eta>0$ and first suppose
$r=\sqrt{u^2+v^2}\le\eta\lambda$.
Each prime $P\le p\le2P$ has an angular density proportional to
$e^{a_p\lambda\cos\theta}$ whose parameter ranges in a fixed compact
subinterval of $(0,\infty)$. Its covariance matrix for
$(\cos\theta,\sin\theta)$ is uniformly positive definite.
The Taylor expansion of its characteristic function in modulus,
uniform in this compact family, bounds each factor by
$e^{-c a_p^2r^2}$. Independently, the fixed prime $p=3$ contributes in
modulus at most $C e^{-c v^2/\lambda}$. Indeed its factor is
$I_0(a_3w)/I_0(a_3\lambda)$; the sectorial expansion bounds its
modulus by $C e^{a_3(\Re w-\lambda)}$, and
$\Re(1/(w+\lambda+iu))\ge c/\lambda$ throughout this disk.
All remaining factors have modulus at most one. Hence
\[
 |\phi(u,v)|\le C\exp(-cM u^2/\lambda^2-cv^2/\lambda)
 \le C\exp(-c'(v_1u^2+v_2v^2)).
\]

On $\eta\lambda\le r\le R\lambda$ each transition-prime factor has
modulus at most $q<1$, uniformly for fixed $R$. In fact its frequency
$a_pr$ ranges in a compact subset of $(0,\infty)$, and equality in
the triangle inequality is impossible on a circle with strictly positive
angular density. For $r\ge R\lambda$ each such factor is at most
$C(r/\lambda)^{-1/2}$. This is the circle Fourier integral bound:
rotate the oscillation to $e^{ia_pr\cos\theta}$, split off intervals of
length $(a_pr)^{-1/2}$ at its two stationary points, and integrate by
parts on the rest. The angular weights and their first derivatives are
uniformly bounded. Choose $R$ large enough that $CR^{-1/2}<1$.
Polar integration of the products, using $M>4$, gives
\[
 \int_{r>\eta\lambda}|\phi(u,v)|\,du\,dv
 \ll \lambda^2e^{-c_R M}.
\]
Thus the characteristic function is integrable. The contribution of this
last domain after rescaling is at most
$O(\sqrt{v_1v_2}\lambda^2e^{-c_R M})=o(1)$.
On the remaining domain, the rescaled integrand has the fixed integrable
Gaussian majorant above; the domain exhausts $\mathbb R^2$.
Fourier inversion and dominated convergence now give
\[
 \sqrt{v_1v_2}\,g(x,0)\longrightarrow
 (2\pi)^{-2}\int_{\mathbb R^2}e^{-(s^2+t^2)/2}\,ds\,dt
 =\frac1{2\pi}.
\]
Substitution into the radial reduction proves the stated formula.
\end{proof}

\section{Harmonic defects on the conditioned prime-zeta edge}
\label{sec:edge-harmonics}

The fixed-harmonic defect of \cite{vm149} measures phase concentration
in a deterministic prime-reciprocal model. We apply the same observable
to the random prime phases at the boundary zero studied in
\cite{vm004}. Keep the notation of Section~\ref{sec:primeedge}. Write
$U_p=e^{i\theta_p}$, $A=\sum_{p\ge3}a_p$, and, for each fixed integer
$m\ge1$, set
\[
 H_m=A-\left|\sum_{p\ge3}a_pU_p^m\right|,\qquad
 B=\frac{\lambda^{1/\sigma-1}}{\log\lambda},\qquad
 C_m=\int_0^\infty y^{-1/\sigma_e}
             \left(1-\frac{I_m(y)}{I_0(y)}\right)\,dy .
\]

\begin{theorem}
On the exact fiber $S=(x,0)$ with the density-conditioned phase law
specified below, $H_m/B\to C_m$ in probability as
$\sigma\uparrow\sigma_e$. This holds jointly for any finite list of
harmonics and also under the conditional law weighted by the
prime-zeta derivative squared. In particular the limit in probability
of $H_m/H_1$ equals $C_m/C_1$, strictly between zero and $m^2$
for $m\ge2$.
\end{theorem}

\begin{proof}
Under $\mathbb P_\lambda$ the angles are independent with density
proportional to $e^{y\cos\theta}$ for $y=\lambda a_p$.
Write $h_m(y)=1-I_m(y)/I_0(y)$,
$D_m^*=\sum a_p(1-\cos m\theta_p)$ and $Y_m=\sum a_p\sin m\theta_p$.
The Fourier integral for $I_m$ gives
$\mathbb E_\lambda D_m^*=\sum a_ph_m(\lambda a_p)$.
The function $h_m$ is continuous and positive, and
$h_m(y)\ll_m\min(1,1/y)$. For the latter estimate use
$1-\cos m\theta\le m^2(1-\cos\theta)$ and the Gaussian bound
$e^{y\cos\theta}/(2\pi I_0(y))\ll\sqrt y\,e^{-cy\theta^2}$
for $|\theta|\le\pi$, $y\ge1$.
Put $P=\lambda^{1/\sigma}$. For $0<\delta<1<K$, partial summation
and the prime number theorem on $\delta P\le p\le KP$ give, after
division by $B$,
\[
 \sigma\int_\delta^K t^{-\sigma}h_m(t^{-\sigma})\,dt+o(1).
\]
Indeed the rescaled prime counting measure
$(\log P/P)\sum_p\delta_{p/P}$ tends to Lebesgue measure on compact
subintervals of $(0,\infty)$, and the integrands converge uniformly
there. Primes below $\delta P$ cost $O_m(\pi(\delta P)/\lambda)$;
the upper tail costs $O(\sum_{p>KP}p^{-\sigma})$.
After division by $B$, their limiting upper bounds are respectively
$O_m(\delta)$ and $O(K^{1-\sigma_0})$ for fixed $\sigma_0>1$
below the working range, by dyadic summation and the prime number
theorem. Send $\delta\downarrow0$, $K\uparrow\infty$, and then use
$y=t^{-\sigma_e}$ to obtain $\mathbb E_\lambda D_m^*/B\to C_m$.
Both endpoints of this integral converge by the stated estimate.

With $M=\#\{p:P\le p\le2P\}$, the same Gaussian bound gives
$\mathop{\rm Var}_\lambda(1-\cos m\theta_p)
\ll_m \min(1,(\lambda a_p)^{-2})$.
Splitting at $P$ as in Section~\ref{sec:primeedge} gives
$\mathop{\rm Var}_\lambda D_m^*\ll_m M/\lambda^2$.
The angular law is even, and
$\sin^2(m\theta)\le m^2\sin^2\theta$ gives
$\mathop{\rm Var}_\lambda Y_m\le m^2v_2=O_m(1/\lambda)$.
These variance upper bounds also apply to either subsum in a partition
of the primes. Since $B\asymp M/\lambda$, they imply concentration
of $D_m^*$ on scale $o(B)$ and $Y_m=O_{\mathbb P}(\lambda^{-1/2})$
before conditioning.

Partition the primes $p\ge3$ into two blocks by alternating in their
increasing order; write $S_j$ for their vector sums and $g_j$ for the
tilted densities. Every block contains a fixed prime and
$M/2+O(1)$ transition primes in $[P,2P]$. The small-disk and
large-frequency majorants in Section~\ref{sec:primeedge}, applied to
each block separately, give
$\|g_j\|_\infty\ll(v_1v_2)^{-1/2}$ by Fourier inversion.
For a block event $E_j$ the continuous density of the joint weighted
measure at $z$ is
$\mathbb E_\lambda[\mathbf1_{E_j}g_{3-j}(z-S_j)]$.
Dividing by the full tilted density $g(z)$ specifies its conditioned
probability at $z$. Equivalently these are the block marginals obtained
by conditioning on shrinking balls about $z$. One may specify the
full fiber law by cylinder events, using convolution with the
continuous remaining-prime density after deleting finitely many
primes. The resulting consistent finite-dimensional measures define
a probability law; uniform convergence of the prime sums puts it on
the stated fiber. Its block marginals agree with the displayed formula
by agreement on cylinders and uniqueness. At $z=(x,0)$ the local
limit theorem gives $g(x,0)\asymp(v_1v_2)^{-1/2}$, hence every
conditioned block event has probability $O(\mathbb P_\lambda(E_j))$.
Apply Chebyshev to the two block deviations from their respective
tilted means, and then the union bound. This preserves
$D_m^*/B\to C_m$ and $Y_m=O_{\mathbb P}(\lambda^{-1/2})$ on the fiber.

Now $A$ stays bounded away from zero and $B\to0$, so $A-D_m^*>A/2$
with probability tending to one. On this event,
\[
 0\le D_m^*-H_m
 =\frac{Y_m^2}{|\sum a_pU_p^m|+A-D_m^*}\ll Y_m^2.
\]
Since $1/\lambda=o(B)$, this proves the assertion for $H_m$.
The tilt is constant on each fiber, so the original and tilted
density-conditioned laws agree there. The fiber identity
$\sum a_p|1-U_p|^2=2(A-x)$ and Cauchy--Schwarz give, with
$U_2=-1$,
\[
 \left|\,\left|\sum_{p\ge2}(\log p)p^{-\sigma}U_p\right|
             -|D_1(\sigma)|\,\right|
 \le \left(2(A-x)\sum_{p\ge3}(\log p)^2p^{-\sigma}\right)^{1/2}.
\]
The bound is uniform on the entire fiber and tends to zero, while
$D_1(\sigma_e)\ne0$. Normalizing the squared-derivative weight
therefore changes probabilities by $o(1)$ uniformly in the event.
Rotation to the aligned frame leaves all $H_m$ invariant.
Finally $1-\cos(m\theta)<m^2(1-\cos\theta)$ away from
$2\pi\mathbb Z$ for $m\ge2$, by the strict triangle inequality for
$\sum_{j=0}^{m-1}e^{ij\theta}$. Positive angular density gives
$0<h_m(y)<m^2h_1(y)$ for $y>0$. Integrating gives the strict ratio
bound; a finite union bound proves the joint assertion.
\end{proof}

\section{The periodic case of Csordas and Escassut's operator problem}
\label{sec:periodic-detector}

Csordas and Escassut's Open Problem~3.3 asks which Laguerre--P\'olya
symbols $\varphi$ of order less than two have the following property:
if a real entire function $f$ in their canonical strip class satisfies
\[
 L[\varphi(\mu D)f](x)>0
 \quad (x\in\mathbb R,\ \mu\in\mathbb R\setminus\{0\}),
 \qquad L[g]=(g')^2-gg'',\qquad D=d/dx,
\]
then $f$ belongs to the Laguerre--P\'olya class \cite{ce}.
We obtain the exact classification on the entire periodic sector of the
original strip class, allowing arbitrary period. The
following proof is independent of any unverified corpus theorem.

\begin{theorem}
Let $\varphi$ be a nonzero real Laguerre--P\'olya function of order
less than two. The following two statements are equivalent.
\begin{enumerate}
\item For every $h>0$ and every real trigonometric polynomial
$f(x)=\sum_{k=-d}^d c_k e^{ikhx}$,
$c_{-k}=\overline{c_k}$, strict positivity of
$L[\varphi(\mu D)f](x)$ for all real $x$ and all real $\mu\ne0$
implies that every zero of $f$ is real.
\item $\varphi(0)\ne0$ and $\varphi$ has infinitely many zeros.
\end{enumerate}
The functions in (1) exhaust the periodic members of the original strip
class. The same classification holds if a single $h>0$ is fixed at the outset.
In particular, condition (2) is necessary in the original problem
with the full strip class.
\end{theorem}

\begin{proof}
First let $f$ belong to the canonical strip class and have period $T>0$.
There are only finitely many zero orbits modulo $T$: representatives
lie in the compact rectangle $0\le\Re z\le T$, $|\Im z|\le A$.
Choose exactly one representative $z_j$ per orbit, with multiplicity
$m_j$. Then $P(w)=\prod_j(w-e^{2\pi i z_j/T})^{m_j}$ has the property
that $P(e^{2\pi iz/T})$ has precisely the zeros of $f$, with their
multiplicities. Both entire functions have finite order at most two
by the canonical strip representation and the exponential-polynomial
form, respectively. Hadamard factorization, using the same genus-two
canonical factors for their common zeros, yields
$f(z)=e^{Q(z)}P(e^{2\pi iz/T})$ for a polynomial $Q$ of degree at
most two. Periodicity forces $Q(z+T)-Q(z)$ to be a fixed integral
multiple of $2\pi i$. Thus $Q(z)=2\pi i n z/T+q_0$ for some integer
$n$. Consequently $f$ is a Laurent polynomial in $e^{2\pi iz/T}$,
and reality gives conjugate Fourier coefficients.

It suffices now to take $h=1$ by rescaling the variable and $\mu$. For $d>0$ take $c_d\ne0$. The Laurent representation
$f(z)=e^{-idz}P(e^{iz})$ has $\deg P=2d$ and $P(0)\ne0$.
Its zeros lie on finitely many horizontal lines and form finitely
many translates of $2\pi\mathbb Z$, with multiplicities. Thus their
nonzero inverse squares are absolutely summable. The function has
exponential type and its real Hadamard factorization puts it in the
strip class in \cite{ce}. Within this class reality of all zeros is
equivalent to membership in the Laguerre--P\'olya class. Constant
test functions never satisfy the strict premise.

Suppose (2) holds. There is no Gaussian factor since the order is
less than two. Writing the real nonzero zeros of $\varphi$ as $r_j$,
with multiplicities, its canonical product gives
\[
 |\varphi(it)|=|\varphi(0)|\prod_j(1+t^2/r_j^2)^{1/2}.
\]
For $0\le k<d$ this implies
\[
 \frac{|\varphi(idt)|}{|\varphi(ikt)|}\longrightarrow\infty
 \qquad(t\longrightarrow+\infty).
\]
Indeed every factor in the ratio is at least one, and for $k>0$
the product of any $M$ factors tends to $(d/k)^M$. For $k=0$
one factor suffices. Put
\[
 G_t(x)=\varphi(tD)f(x)
       =\sum_{k=-d}^d c_k\varphi(ikt)e^{ikx}.
\]
After division by $2|c_d\varphi(idt)|$, $G_t$ differs uniformly from
$\cos(dx+\arg(c_d\varphi(idt)))$ by $o(1)$ as $t\to\infty$.
Hence at the $2d$ consecutive extrema of this cosine in one period
the signs alternate for all sufficiently large $t$. There are
$2d$ distinct real zeros modulo the period. Since the polynomial
$P_t(w)=\sum_{k=-d}^d c_k\varphi(ikt)w^{k+d}$ has degree $2d$
and a nonzero constant coefficient, these are all its roots under
$w=e^{ix}$ and all are simple.

If $f$ had a nonreal zero, $P_0$ would have an off-circle root,
since $\varphi(0)\ne0$. Fix $T$ sufficiently large as above.
The set of $t\in[0,T]$ at which all roots of $P_t$ lie on the unit
circle is closed by continuity of polynomial roots. Its minimum
$t_*$ is positive. At $t_*$ at least one root must be multiple:
if all were simple, all $2d$ distinct zeros of the real function
$G_{t_*}(x)$ in a period would persist as real zeros for all nearby
$t$, by sign changes in disjoint small intervals. This contradicts
minimality. A multiple root on the unit circle yields a real $x_*$
with $G_{t_*}(x_*)=G'_{t_*}(x_*)=0$, so $L[G_{t_*}](x_*)=0$.

For the converse, if $\varphi(0)=0$, use $f(x)=2+\sin x$.
It has nonreal zeros, whereas $G_t$ is a sinusoid of nonzero
amplitude $|\varphi(it)|$ for every $t\ne0$. Indeed the nonzero
zeros of a Laguerre--P\'olya symbol are real. Its Laguerre expression
is the amplitude squared.

It remains to consider $\varphi(0)\ne0$ with only $m<\infty$ zeros.
Then $\varphi(z)=e^{bz}P(z)$ for a real $b$ and a real polynomial
$P$ of degree $m$ with only real nonzero zeros. Thus
$|\varphi(2it)|/|\varphi(it)|\le2^m$ for every real $t$, by
comparison of the factors associated to each zero. Set
$\varepsilon=2^{-m}/20$ and
\[
 f(x)=\sin x+\varepsilon\sin(2x)
     =\sin x(1+2\varepsilon\cos x).
\]
The second factor has nonreal zeros since $1/(2\varepsilon)>1$.
For each $t$, division of $G_t$ by $|\varphi(it)|$ gives
$g(x)=\sin(x+\alpha)+r\sin(2x+\beta)$ for phases depending on
$t$ and $0\le r\le 1/20$. With $u=x+\alpha$, $v=2x+\beta$,
\[
 L[g](x)=1+r(4\cos u\cos v+5\sin u\sin v)+4r^2
       \ge 1-5r\ge 3/4.
\]
The absolute value of the mixed term in parentheses is at most five
by the Cauchy--Schwarz inequality. Thus $f$ violates the detector
implication. Rescaling these examples proves the assertion for
each fixed period.
\end{proof}

The cosine result and the obstruction at $\varphi(0)=0$ also appear
in the corpus nodal criterion \cite{vm103}. The Gaussian collision
argument studies a symbol of order two \cite{vm126}. Here its
multiple-real-zero test has a finite-dimensional counterpart: the
Laurent polynomial has fixed degree and its unit-circle roots can be
counted in one period. Dominance of the highest Fourier mode supplies
the all-real endpoint for every admissible symbol in (2). A general
strip function has no assumed highest Fourier mode. The proof
therefore leaves sufficiency of (2) for the full strip class open.
The characterization concerns this restricted domain. Its historical
priority has not been established.

\subsection{An infinite spectrum with an attained endpoint}
\label{sec:purepoint}

\begin{theorem}
The same symbol classification holds for all real functions
\[
 f(z)=\sum_{\lambda\in E}c_{\lambda}e^{i\lambda z},\qquad
 E=-E\subset[-d,d],\quad 0<d\in E,\quad
 c_{-\lambda}=\overline{c_{\lambda}},\quad
 c_d\ne0,\quad \sum_{\lambda\in E}|c_\lambda|<\infty .
\]
Here $E$ may be countably infinite. Frequencies may be
incommensurate and may accumulate at either endpoint.
\end{theorem}

\begin{proof}
Necessity follows from the two periodic counterexamples. Assume that
$\varphi(0)\ne0$ and that $\varphi$ has infinitely many zeros.
Put $G_t(z)=\sum_{\lambda} c_\lambda\varphi(i\lambda t)e^{i\lambda z}$.
The series and its spatial derivatives converge locally uniformly in
the two complex variables $(t,z)$. Indeed $|\lambda|\le d$ and
the coefficients are absolutely summable. Applying the entire power
series of $\varphi$ to $tD$ termwise is justified by
$\sum_n |[z^n]\varphi|\,(d|t|)^n<\infty$ and the same summability.
For real $t$ the function $G_t$ is real on the real axis. The two
edge coefficients $c_{\pm d}\varphi(\pm idt)$ are nonzero. For each
fixed real $t$, on $\Im z=y\to+\infty$ the ratio of the sum of the
remaining terms to the $-d$ term tends uniformly to zero, by
dominated convergence of the series of absolute values. The $+d$
term gives the corresponding lower-half-plane bound. Thus $G_t$
has all zeros in some bounded horizontal strip. It is of exponential
type; Jensen's formula and Hadamard factorization therefore place it
in the canonical strip class.

Suppose $f$ has a nonreal zero and that all the strict Laguerre
inequalities hold. Since $G_0=\varphi(0)f$, Rouch\'e's theorem
gives a $t_0>0$ such that $g=G_{t_0}$ has a nonreal zero. All real
zeros of $G_t$ are simple at every positive $t$. For $|\lambda|<d$,
the canonical product in the preceding proof gives
\[
 0\le\frac{|\varphi(i\lambda t)|}{|\varphi(idt)|}\le 1,
 \qquad
 \frac{|\varphi(i\lambda t)|}{|\varphi(idt)|}\longrightarrow0
 \quad(t\to+\infty).
\]
Dominated convergence now shows that at some $T>t_0$ the real
sinusoid formed from the extreme pair of $G_T$ has larger amplitude
than the sum of the amplitudes of all other modes. Its alternating
extrema force at least $2Rd/\pi-O(1)$ distinct real zeros of
$G_T$ in every interval of length $2R$.

Choose $Y$ large enough for the relative errors in both extreme-term
approximations of $g$ to be bounded by $1/2$ for $|\Im z|\ge Y$.
Choose $x_0\in\mathbb R$ such that the vertical segment
$x_0+i[-Y,Y]$ contains no zero of $g$ and $G_T(x_0)\ne0$.
Only countably many choices are excluded. There exist $R_j\to\infty$
with $e^{i\lambda R_j}\to1$ for every $\lambda\in E$. To see this,
take the closure of $s\mapsto(e^{i\lambda s})_{\lambda\in E}$ in
the compact product of unit circles. For any symmetric identity
neighborhood $U$ in this closure, finitely many translates by image
points cover the closure. Hence every real $s$ lies within bounded
distance (depending on $U$) of a return time to $U$. Enumerate
$E$ and take successive neighborhoods restricting the first $j$
coordinates to distance $1/j$ of the identity; the corresponding
times can be chosen arbitrarily large.

For either sign the endpoint functions $t\mapsto G_t(x_0\pm R_j)$
converge to $t\mapsto G_t(x_0)$ uniformly on compact sets of complex
$t$, by dominated convergence of the coefficients. The limit function
is nonzero at $t_0$. Take a closed disk in the complex $t$-plane
having $[t_0,T]$ in its interior and no boundary zero of the limit.
Rouch\'e's theorem gives a uniform bound on the number of real
endpoint zeros in $[t_0,T]$. Moreover both endpoints have nonzero
values at $t_0$ and $T$ for all sufficiently large $j$.
On each fixed time--space rectangle with spatial base
$I_j=[x_0-R_j,x_0+R_j]$, the number of real zeros can change only at
these endpoint events. Every interior zero is simple in space and
locally persists, by the implicit function theorem and compactness.
At an endpoint event the spatially simple zero can change the count
by at most one per endpoint. Thus the real root counts at $t_0$ and
$T$ in $I_j$ differ by $O(1)$ independently of $j$.

The total number of zeros of $g$, counted with multiplicity over
$\Re z\in I_j$, is $2R_jd/\pi+O(1)$. Indeed the bottom boundary of
the rectangle $I_j+i[-Y,Y]$, oriented to the right, contributes
$2R_jd+O(1)$ to the net change of argument by the choice of $Y$;
the top side, oriented to the left, contributes the same amount.
On the vertical sides $g(x_0\pm R_j+iy)$ and their $y$-derivatives
converge uniformly to those on $x_0+i[-Y,Y]$, again by absolute
summability and bounded frequencies. The reference segment is
zero-free, so the two vertical changes of argument are bounded,
for example by integrating the absolute logarithmic derivative.
The argument principle proves the count. Together these estimates
give at most $O(1)$ nonreal zeros in the expanding intervals $I_j$.

On the other hand, a single nonreal zero of $g$ produces linearly
many nonreal zeros. Choose a small closed disk about it disjoint from
the real axis and zero-free on its boundary. Absolute convergence on
the disk first controls the tail of the Fourier series uniformly over
all real translates. Requiring the finitely many remaining phases
$e^{i\lambda s}$ to be sufficiently close to $1$ therefore makes
$g(z+s)$ uniformly close enough to $g(z)$ for Rouch\'e's theorem.
The compact group argument above supplies a relatively dense set of
such $s$. Taking shifts in sufficiently widely spaced intervals gives
a number $\gg R_j$ of disjoint translated disks with real parts in
$I_j$, each containing a nonreal zero. This contradicts the preceding
bound and proves sufficiency.
\end{proof}

The proof counts crossings on two selected vertical segments, as
in the nodal method of \cite{vm103}. Here simultaneous recurrence
makes both boundary functions converge, while the atoms at $\pm d$
give the exact zero-count slope. Finite exponential sums and all
periodic members of the strip class are included. A pure point
spectrum accumulating at an unattained endpoint, a continuous
spectrum, and a general canonical strip product need not supply
the dominating sinusoid used in this proof.

\section{A reciprocal-zero trace test for trigonometric polynomials}
\label{sec:trace-circle}

The zeros of a real trigonometric polynomial in one complex period
strip can be counted by exact operations on its Fourier coefficients.
This is the Hermite trace-form argument with the unit-circle reflection
in place of ordinary conjugation. It gives a finite diagnostic for
periodic multiplier families, and uses the reciprocal-root traces that
also occur in \cite{vm002}. We give the derivation to include multiple
roots and changes of degree.

Suppose
\[
 F_t(z)=\sum_{j=-d}^{d}c_j(t)e^{ijz},\qquad
 c_{-j}(t)=\overline{c_j(t)},\qquad c_d(t)\ne0.
\]
Put $P_t(w)=\sum_{j=-d}^{d}c_j(t)w^{j+d}$ and $D=2d$.
The roots of $P_t$ occur with their reflected partners
$\iota(w)=1/\overline w$, with equal multiplicity.
The map $w=e^{iz}$ identifies unit-circle roots with real zeros of $F_t$
modulo $2\pi$; roots strictly inside the circle correspond to upper
half-plane zeros. Write $s_h=\sum_{P_t(\alpha)=0}\alpha^h$, counting
multiplicity, and form the Hermitian Toeplitz matrix
\[
 H_t=(s_{j-i})_{0\le i,j<D},\qquad s_0=D,\qquad
 s_{-h}=\overline{s_h}.
\]
Only $s_1,\ldots,s_{D-1}$ are required. Newton's identities compute them
from the coefficients of the monic normalization of $P_t$.

\begin{proposition}
If $a$ is the number of distinct roots on the circle and $b$ is the
number of distinct two-element reflection orbits, then the positive,
negative and zero inertia indices of $H_t$ are
\[
 (a+b,\ b,\ D-a-2b).
\]
In particular $H_t$ is positive semidefinite exactly when all zeros
of $F_t$ are real modulo period, and positive definite exactly when
they are also simple. Singularity, whether on or off the circle,
detects repeated roots.
\end{proposition}

\begin{proof}
For $q(w)=\sum_{j=0}^{D-1}x_jw^j$ the quadratic form is
\[
 x^*H_tx=\sum_{\alpha}m_\alpha
 \overline{q(\iota(\alpha))}q(\alpha),
\]
where the sum is over distinct roots. Evaluation on these roots is
onto by Lagrange interpolation. A fixed root contributes one positive
square and a two-element orbit contributes the two-dimensional form
$2m\Re(\overline{u}v)$ of inertia $(1,1)$. The kernel of evaluation
accounts for the remaining dimensions.
\end{proof}

For continuous coefficients of constant degree the inertia is locally
constant at every nonsingular time. A singular time is a candidate for
a change in the number of nonreal pairs. The negative indices on the
two adjacent simple-root intervals give the actual change. A singular
matrix alone cannot locate the collision on the circle. If multiplicity
counts are wanted at a singular time, use the squarefree layers
$P_t=\mathrm{const}\prod_{m\ge1}Q_m^m$, where $Q_m$ is monic and has
as its simple roots the roots of exact multiplicity $m$. Each $Q_m$
retains the reflection symmetry, so $\sum_m m\,n_-(H(Q_m))$ counts
upper half-plane zeros modulo period with multiplicity. If the leading
frequency vanishes, first reduce to the actual Fourier degree.
A nonzero constant has no zeros; an identically zero slice must be
excluded.

For example the backward Gaussian heat flow of $\cos z-c$, $c>1$, is
$e^t\cos z-c$. The monic circle polynomial is
$w^2-2ce^{-t}w+1$ and its trace matrix has eigenvalues
$2\pm2ce^{-t}$. For $0\le t<\log c$ the negative index is one.
At $t=\log c$ there is a double root at $w=1$ and hence a real collision
at $z=0$ modulo $2\pi$. For $t>\log c$ the trace matrix is positive
definite. More generally, algebraic Fourier coefficients permit exact
sign determination by algebraic arithmetic whenever the multiplier
values are algebraic. This computation concerns finite Fourier spectra.

\section{Compatibility of decoder laws in arbitrary magmas}
\label{sec:decoder}
For a nonempty magma write $L_y(u)=y*u$ and $R_y(u)=u*y$. With
$C_y(x,z)=x*(z*y)$, manuscript~\cite{vm085} studies the laws
\[
 E_k:\quad L_y^k C_y(x,z)=x,\qquad
 G_k:\quad L_y^k C_y(x,z)=z\qquad(k\ge1).
\]
The notation for $G_k$ exchanges the dummy variables in the source.
The next assertion places no finiteness bound on a common model.

\begin{theorem}
Any two distinct laws in $\{E_k,G_k:k\ge1\}$ have only the one-element
magma as a common model. For each $k,l\ge1$, the varieties defined by
$E_k$ and $G_l$ have a binary term selecting the first input in the
former and the second in the latter.
\end{theorem}
\begin{proof}
Either law makes every $L_y$ surjective, because the output of its
$k$-fold iterate is an arbitrary element. Every $R_w$ is injective
as well. For $G_k$ we have $L_w^k L_x R_w=\mathrm{id}$ for fixed
$x$. For $E_k$, choose $s,t$ with $w=s*t$, which is possible by
surjectivity of $L_s$ for any chosen $s$. Then
$L_t^k R_w=\mathrm{id}$.

For $E_k$ and $G_l$, applying the longer decoder to $C_y(x,z)$ yields
$z=L_y^{l-k}x$ when $l\ge k$, and $x=L_y^{k-l}z$ when $k>l$.
Either identity for independently varying $x,z$ forces a singleton.
For two laws of the same family with depths $k<l$, their common
output ranges over the magma and satisfies $L_y^{l-k}u=u$ for every
$y,u$. Thus all $L_y$ are injective. In the $E$ case, fixing $x,y$
in $L_y^k L_x(z*y)=x$ now makes $z*y$ independent of $z$.
The injectivity of $R_y$ forces a singleton. In the $G$ case, fixing
$z,y$ in $L_y^k R_{z*y}(x)=z$ gives a constant injective function
of $x$, again forcing a singleton.

Set $D_j(y,x,z)=L_y^j C_y(x,z)$ and, for fixed $y$, put
$S_{j,y}(w)=L_y^{j-1}C_y(w,w)$. In either depth-$j$ variety
$L_yS_{j,y}=\mathrm{id}$, hence $L_y^n S_{j,y}^n=\mathrm{id}$ for
every $n\ge0$. Use $y=u$ in the following binary term:
\[
 s(u,v)=
 \begin{cases}
 D_k(y,u,S_{l,y}^{k-l}(v)),& k\ge l,\\
 D_l(y,S_{k,y}^{l-k}(u),v),& k<l.
 \end{cases}
\]
In the first case $E_k$ gives $u$ immediately and $G_l$ gives
$L_y^{k-l}S_{l,y}^{k-l}(v)=v$. In the second case $G_l$ gives $v$
immediately and $E_k$ gives $L_y^{l-k}S_{k,y}^{l-k}(u)=u$.
\end{proof}

Thus the $E_k$ and $G_l$ varieties are independent in the standard
binary-term sense. The proof uses the equations alone. The
nontrivial infinite models for individual depths $k\ge3$ are supplied
by the decoder-tree construction reported in \cite{vm085}.

\section*{Where the corpus methods enter}
\addcontentsline{toc}{section}{Where the corpus methods enter}
In the Diophantine argument the normalized remainder of \cite{vm090}
turns a specified word of binary choices into one terminal equation.
A residue condition propagates that word to a whole progression.
This use of congruence propagation is elementary; the two-state
$p$-adic operator for Sun's sums in \cite{vm151} and the harmonic
mass lemma for Mills prime trees in \cite{vm141} are not hypotheses.

Two transfers use the same observation in several parts of one proof.
Section~\ref{sec:edge-harmonics} applies the rotation-invariant
deficit of \cite{vm149} on the phase-locked fiber of \cite{vm004};
the local limit and the split-prime estimate keep that statistic
under control after conditioning. In Section~\ref{sec:periodic-detector}
a multiple real zero witnesses a collision as in \cite{vm126};
Section~\ref{sec:trace-circle} expresses its finite-frequency test
in reciprocal-root power sums, a representation also used in
\cite{vm002}. The invariance underlying a diagnostic has to be checked:
for example, the trace matrix is singular at every multiple root,
and only its inertia distinguishes pairs away from the circle.
This limitation is proved by the trace form itself.

In the corner problem we use the threshold-summation strategy of
\cite{vm161} with total degree as the threshold. The finite monomial
theorem is classical \cite{ggr}. Its equality cases on all nested
triangles are forced by a layer recurrence. In Section~\ref{sec:decoder}
the two recovery rules of \cite{vm085} furnish an actual binary selector.
The arithmetic shear in \cite{vm020} concerns a different obstruction
and is not used to infer any statement about Euler's constant.
These distinctions identify the imported theorems, adapted measurements,
and new steps separately.

\begin{thebibliography}{99}\small
\bibitem{tuz} S.~Tengely, M.~Ulas and J.~Zygad\l o,
\emph{On a Diophantine equation of Erd\H{o}s and Graham},
J. Number Theory \textbf{217} (2020), 445--459.
\href{https://doi.org/10.1016/j.jnt.2020.05.006}{doi:10.1016/j.jnt.2020.05.006}.
\bibitem{bl} P.~B. Borwein and T.~A. Loring,
\emph{Some questions of Erd\H{o}s and Graham on numbers of the form
$\sum g_n/2^{g_n}$}, Math. Comp. \textbf{54} (1990), 377--394.
\bibitem{ce} G.~Csordas and A.~Escassut,
\emph{The Laguerre inequality and the distribution of zeros of entire functions},
Ann. Math. Blaise Pascal \textbf{12} (2005), 331--345.
\href{https://doi.org/10.5802/ambp.210}{doi:10.5802/ambp.210}.
\bibitem{ggr} A.~V. Geramita, D.~Gregory and L.~Roberts,
\emph{Monomial ideals and points in projective space},
J. Pure Appl. Algebra \textbf{40} (1986), 33--62, Theorem 5.4(1).
\href{https://doi.org/10.1016/0022-4049(86)90029-0}{doi:10.1016/0022-4049(86)90029-0}.
\bibitem{dlmf} NIST Digital Library of Mathematical Functions,
\href{https://dlmf.nist.gov/10.32}{\S10.32}
and \href{https://dlmf.nist.gov/10.40}{\S10.40}, modified Bessel functions.
\bibitem{vmrepo} QEEDDD, \emph{Verifiable Mathematical Reasoning Corpus},
\href{https://github.com/QEEDDD/Verifiable-Mathematical-Reasoning-Corpus/tree/5887cdea4898f0a3aa5bf7a8e0829697bd2b0f4b}{revision 5887cdea},
CC BY 4.0. Entries below give manuscript numbers in this revision.
\bibitem{vm002} VMRC [002], \href{https://github.com/QEEDDD/Verifiable-Mathematical-Reasoning-Corpus/blob/5887cdea4898f0a3aa5bf7a8e0829697bd2b0f4b/\%5B002\%5D\%20The\%20Polyexponential\%20Trace-Tower\%20Program\%20\%28Moving\%20\%20Levels\%2C\%20Faber\%20and\%20Inverse\%20Monodromy\%2C\%20Positive\%20Spectra\%2C\%20and\%20Arithmetic\%20Barriers\%29.tex}{The polyexponential trace-tower program}, research manuscript in \cite{vmrepo}.
\bibitem{vm004} VMRC [004], \href{https://github.com/QEEDDD/Verifiable-Mathematical-Reasoning-Corpus/blob/5887cdea4898f0a3aa5bf7a8e0829697bd2b0f4b/\%5B004\%5D\%20Prime-Zeta\%20Divisor\%20Tomography\%20Local\%20Zero\%20Processes\%2C\%20Frozen-Edge\%20Rigidity\%2C\%20and\%20Derivative\%20Caustics.tex}{Prime-zeta divisor tomography}, research manuscript in \cite{vmrepo}.
\bibitem{vm020} VMRC [020], \href{https://github.com/QEEDDD/Verifiable-Mathematical-Reasoning-Corpus/blob/5887cdea4898f0a3aa5bf7a8e0829697bd2b0f4b/\%5B020\%5D\%20The\%20Two-World\%20Search\%20for\%20Euler\%E2\%80\%99s\%20Constant\%20\%28Shear\%20invariance\%2C\%20collision\%20worlds\%2C\%20forgeries\%2C\%20and\%20the\%20map\%20of\%20possible\%20worlds\%29.tex}{The two-world search for Euler's constant}, research manuscript in \cite{vmrepo}.
\bibitem{vm085} VMRC [085], \href{https://github.com/QEEDDD/Verifiable-Mathematical-Reasoning-Corpus/blob/5887cdea4898f0a3aa5bf7a8e0829697bd2b0f4b/\%5B085\%5D\%20decoder_trees_5093_5107.tex}{Decoder trees and Kisielewicz's identity}, research manuscript in \cite{vmrepo}.
\bibitem{vm090} VMRC [090], \href{https://github.com/QEEDDD/Verifiable-Mathematical-Reasoning-Corpus/blob/5887cdea4898f0a3aa5bf7a8e0829697bd2b0f4b/\%5B090\%5D\%20erdos_graham_tent.tex}{The Erd\H{o}s--Graham tent recurrence and dyadic representations}, research manuscript in \cite{vmrepo}.
\bibitem{vm103} VMRC [103], \href{https://github.com/QEEDDD/Verifiable-Mathematical-Reasoning-Corpus/blob/5887cdea4898f0a3aa5bf7a8e0829697bd2b0f4b/\%5B103\%5D\%20A_Nodal_Criterion_for_Real_Zeros.tex}{A nodal criterion for real zeros of entire functions in a strip}, research manuscript in \cite{vmrepo}.
\bibitem{vm126} VMRC [126], \href{https://github.com/QEEDDD/Verifiable-Mathematical-Reasoning-Corpus/blob/5887cdea4898f0a3aa5bf7a8e0829697bd2b0f4b/\%5B126\%5D\%20Real_Collisions_in_the_Gaussian_Heat_Flow.tex}{Real collisions of zeros under the Gaussian backward heat flow}, research manuscript in \cite{vmrepo}.
\bibitem{vm140} VMRC [140], \href{https://github.com/QEEDDD/Verifiable-Mathematical-Reasoning-Corpus/blob/5887cdea4898f0a3aa5bf7a8e0829697bd2b0f4b/\%5B140\%5D\%20ErdosGraham_Density_Parity_Packings.tex}{Erd\H{o}s--Graham density and parity packings}, research manuscript in \cite{vmrepo}.
\bibitem{vm141} VMRC [141], \href{https://github.com/QEEDDD/Verifiable-Mathematical-Reasoning-Corpus/blob/5887cdea4898f0a3aa5bf7a8e0829697bd2b0f4b/\%5B141\%5D\%20harmonic_branching_mills.tex}{Harmonic branching mass and Mills-type prime-representing functions}, research manuscript in \cite{vmrepo}.
\bibitem{vm149} VMRC [149], \href{https://github.com/QEEDDD/Verifiable-Mathematical-Reasoning-Corpus/blob/5887cdea4898f0a3aa5bf7a8e0829697bd2b0f4b/\%5B149\%5D\%20Prime_Circles_Zeta_Unified_Synthesis\%20\%281\%29.tex}{Prime reciprocal circles and the Riemann zeta function}, research manuscript in \cite{vmrepo}.
\bibitem{vm151} VMRC [151], \href{https://github.com/QEEDDD/Verifiable-Mathematical-Reasoning-Corpus/blob/5887cdea4898f0a3aa5bf7a8e0829697bd2b0f4b/\%5B151\%5D\%20Sun_A179524_Complete_Divisibility_Proof.tex}{A proof of Sun's divisibility conjecture for A179524}, research manuscript in \cite{vmrepo}.
\bibitem{vm161} VMRC [161], \href{https://github.com/QEEDDD/Verifiable-Mathematical-Reasoning-Corpus/blob/5887cdea4898f0a3aa5bf7a8e0829697bd2b0f4b/\%5B161\%5D\%20Randic_Index_and_Vertex_Radius.tex}{A proof of the Caporossi--Hansen conjecture on the Randi\'c index and graph radius}, research manuscript in \cite{vmrepo}.
\end{thebibliography}
\end{document}
